% Figure: one observed failure, four candidate causes, four interventions (public evaluation puzzle 62593bfd).
% Grids: figures/make_same_failure_grids.py (puzzle file and CORA-CAST's recorded decisions); text: the cited records.
\begin{figure}[t]
\centering
\begin{tikzpicture}[
  lab/.style={font=\scriptsize, align=center, inner sep=1pt},
  arr/.style={-{Latex[length=1.5mm]}, draw=black!55, semithick},
  cause/.style={draw=black!60, rounded corners=2pt, fill=black!3, align=left, font=\scriptsize, inner sep=3pt,
    text width=3.62cm, minimum height=3.05cm},
  act/.style={draw=tested!85!black, rounded corners=2pt, fill=tested!9, align=left, font=\scriptsize, inner sep=3pt,
    text width=3.62cm, minimum height=1.05cm},
  obs/.style={draw=pending!85!black, rounded corners=2pt, fill=pending!6, align=center, font=\footnotesize,
    inner sep=4pt, text width=15.6cm}]
% the puzzle: two demonstrations and the first test input with CORA's recorded answer
\node[inner sep=0] (d1i) at (0,0) {\includegraphics[scale=1.0]{sf_demo1_in.pdf}};
\node[inner sep=0, right=0.42cm of d1i] (d1o) {\includegraphics[scale=1.0]{sf_demo1_out.pdf}};
\node[inner sep=0, right=0.85cm of d1o] (d2i) {\includegraphics[scale=1.0]{sf_demo2_in.pdf}};
\node[inner sep=0, right=0.42cm of d2i] (d2o) {\includegraphics[scale=1.0]{sf_demo2_out.pdf}};
\node[inner sep=0, right=0.85cm of d2o] (ti) {\includegraphics[scale=1.0]{sf_test1_in.pdf}};
\node[inner sep=0, right=0.42cm of ti] (to) {\includegraphics[scale=1.0]{sf_test1_cora.pdf}};
\draw[arr] (d1i) -- (d1o); \draw[arr] (d2i) -- (d2o); \draw[arr] (ti) -- (to);
\node[lab, above=1pt of $(d1i.north)!0.5!(d1o.north)$] {demonstration 1 (background colour 0)};
\node[lab, above=1pt of $(d2i.north)!0.5!(d2o.north)$] {demonstration 2 (background colour 5)};
\node[lab, above=1pt of $(ti.north)!0.5!(to.north)$] {test input 1 and CORA's answer (exact)};
% the one observation
\node[obs, below=0.32cm of d2o.south west, anchor=north] (obs) {\textbf{One observation, in five states of the
  system:} no program that reproduces both demonstrations and answers the test. (CORA's original engine: its fit was
  wrong; the original engine with the background fixed; CORA-CAST with one fixed view, with every guard set to true,
  and with its plain typed language: no fit.)};
% four candidate causes, each with the evidence that identified it
\node[cause, below=0.42cm of obs.south west, anchor=north west] (c1) {\textbf{Missing representation.} Six of the
  seven original views took colour 0 as background. The needed rule fits demonstration 1 in three views and
  demonstration 2 only in a 4-connected view that splits demonstration 1's diagonal pieces: no view fits both. One
  fixed view (single-colour, 4-connected): no fit.};
\node[cause, right=0.25cm of c1] (c2) {\textbf{Search failure.} With each grid's majority colour as background the
  rule is expressible and fits both demonstrations, but the original enumerator ranks its two rules 1,332nd and 84th of
  4,320, and no program containing both is among its first 200,000 candidates.};
\node[cause, right=0.25cm of c2] (c3) {\textbf{Wrong binding.} The same views and moves with every guard set to
  true, so that every object takes the same action: no fit, the search exhausted.};
\node[cause, right=0.25cm of c3] (c4) {\textbf{Missing operation.} CORA-CAST's plain typed language, which has no
  guarded move by destination: no fit.};
\foreach \c in {c1,c2,c3,c4} {\draw[arr] (obs.south -| \c.north) -- (\c.north);}
% the intervention each cause calls for
\node[act, below=0.3cm of c1] (a1) {\textbf{Change the representation:} the majority colour as background;
  single-colour 8-connected objects.};
\node[act, below=0.3cm of c2] (a2) {\textbf{Construct instead of enumerating:} build the rule from each object's
  fate in the demonstrations.};
\node[act, below=0.3cm of c3] (a3) {\textbf{Revise the binding:} partition the objects by a feature, whether another
  object has the same size.};
\node[act, below=0.3cm of c4] (a4) {\textbf{Build the operation:} sequential moves until blocked and to the border.};
\foreach \i in {1,2,3,4} {\draw[arr] (c\i) -- (a\i);}
% the resolution
\node[obs, draw=tested!85!black, fill=tested!5, below=0.38cm of $(a2.south)!0.5!(a3.south)$, anchor=north] (res)
  {\textbf{With all four:} \texttt{(rewrite view=mono8 [size\_unique == False -> MOVE up until\_blocked] [TRUE -> MOVE
  down to\_border])}: 8 fitting programs with one behaviour; both test outputs exact.};
\foreach \i in {1,2,3,4} {\draw[arr] (a\i.south) -- (a\i.south |- res.north);}
\end{tikzpicture}
\caption[]{Why the same observed failure can need different interventions: a real case from CORA's records. In
public evaluation puzzle 62593bfd of ARC-AGI-2, objects whose size another object shares move up until blocked and the
others move down to the border. In every state of the system the outcome was the same, no correct program, and the failure
record (no fit, or a fit that is wrong on the test) did not say why; the four causes were told apart only by different
controls: a change of view, a ranking analysis, a guard ablation and an operator ablation. The puzzle was inspected while building and shaped the move constructor, so this is
development evidence, not accuracy. Grids are drawn from the puzzle file and from CORA-CAST's recorded first attempt by
\protect\rec{paper/v2_neural_cora/figures/make_same_failure_grids.py}; the controls' exhausted searches (0 fits each)
are in \protect\rec{experiments/cora_cast_20261026/runs/chain_v2/m2/ARC_eval_dev/decisions.jsonl} and the plain typed
language's in \protect\rec{experiments/cora_cast_20261026/M2_REWRITE_REPORT.md}; the original engine's views and
enumeration ranks are the independent review's analysis (\protect\rec{records/CORA_SYNTHESIS_REVIEW_20261027.md},
sections 2.3 and 3). Puzzle data: ARC Prize Foundation.}
\label{fig:samefailure}
\end{figure}
